\documentclass[12pt]{article}

\usepackage{sbc-template}
\usepackage{graphicx,url}
\usepackage[utf8]{inputenc}
\usepackage[brazil]{babel}
\usepackage{tikz}
\usetikzlibrary{arrows.meta,positioning,shapes.geometric,fit}
\usepackage{booktabs}
\usepackage{array}

\sloppy

\title{Engenharia Reversa de um Sistema Web em Django Sem Documentação Formal: Um Estudo de Caso}

\author{Rafael Zottesso\inst{1}}

\address{
  Instituto Federal do Paraná (IFPR)
  \email{rafael.zottesso@ifpr.edu.br}
}

\begin{document}

\maketitle

\begin{abstract}
  This paper presents a case study on the application of reverse engineering
  techniques to a real Django web system that lacked formal documentation.
  Static code analysis was used to recover functional and non-functional
  requirements, business rules, UML models, a simplified data model, and
  information about the existing automated tests and query optimization
  practices. The results show that a substantial part of the architectural
  and functional documentation of a small/medium web system can be
  reconstructed from source code evidence, while also revealing gaps such as
  the complete absence of automated tests and the lack of empirical
  performance measurements.
\end{abstract}

\begin{resumo}
  Este artigo apresenta um estudo de caso sobre a aplicação de técnicas de
  engenharia reversa em um sistema web real, desenvolvido em Django, que não
  possuía documentação formal. A partir da análise estática do código-fonte,
  foram recuperados requisitos funcionais e não funcionais, regras de
  negócio, modelos UML, um modelo de dados simplificado, além de informações
  sobre os testes automatizados existentes e as práticas de otimização de
  consultas ao banco de dados. Os resultados mostram que parte relevante da
  documentação arquitetural e funcional de um sistema web de pequeno/médio
  porte pode ser reconstruída a partir de evidências do código, ao mesmo
  tempo em que a análise revela lacunas como a ausência total de testes
  automatizados e a falta de medições empíricas de desempenho.
\end{resumo}

\section{Introdução}

Sistemas de software desenvolvidos em contextos acadêmicos, de pequenas
equipes ou de projetos evolutivos frequentemente carecem de documentação
formal atualizada. Requisitos, regras de negócio e decisões arquiteturais
tendem a ficar implícitos no código-fonte, dificultando a manutenção, a
transferência de conhecimento entre desenvolvedores e a avaliação da
qualidade do sistema \cite{chikofsky1990,corbi1989}. Esse cenário é comum em
sistemas web construídos com frameworks produtivos como o Django, nos quais a
velocidade de desenvolvimento tende a ser priorizada em detrimento do
registro formal de requisitos e arquitetura.

A engenharia reversa de software é definida como o processo de analisar um
sistema para identificar seus componentes e relacionamentos, e para criar
representações do sistema em outra forma ou em um nível mais alto de
abstração \cite{chikofsky1990}. Esse processo é particularmente relevante
quando a documentação original é inexistente ou está desatualizada, pois
permite recuperar informações de projeto (\emph{design recovery}) que, de
outra forma, exigiriam entrevistas extensas com os desenvolvedores originais
ou a reconstrução manual e sujeita a erros.

Este trabalho é motivado pela constatação prática de que um sistema web em
produção -- um sistema de gestão de campeonatos de eSports acadêmicos,
implementado em Django -- não possuía nenhum artefato de documentação de
requisitos, arquitetura ou modelagem, além de não possuir testes
automatizados implementados. Diante desse cenário, buscou-se avaliar até que
ponto técnicas de engenharia reversa baseadas em análise estática de código
seriam suficientes para reconstruir uma documentação técnica minimamente
completa, evidenciando também os limites dessa abordagem.

A questão de pesquisa que orienta este trabalho é: \emph{como a engenharia
reversa pode auxiliar na recuperação da documentação arquitetural e
funcional de um sistema web desenvolvido em Django que não possui
documentação formal?}

O objetivo geral deste trabalho é aplicar um processo de engenharia reversa
a um sistema Django real, recuperando requisitos, regras de negócio,
arquitetura, modelos UML, modelo de dados, informações de testes e aspectos
de desempenho, e discutir os resultados obtidos como um estudo de caso. Os
objetivos específicos são: (i) descrever um processo reprodutível de
recuperação de documentação a partir de código-fonte; (ii) aplicar esse
processo a um sistema real; (iii) analisar a qualidade dos testes e das
otimizações de consulta existentes; e (iv) discutir os benefícios e as
limitações da abordagem adotada.

A contribuição deste trabalho é apresentar, de forma sistemática, um relato
de experiência de engenharia reversa aplicada a um sistema Django real,
incluindo os artefatos recuperados (requisitos, regras de negócio, diagramas
UML e matriz de rastreabilidade) e uma discussão crítica sobre o que pôde e o
que não pôde ser recuperado apenas a partir do código-fonte.

O restante deste artigo está organizado da seguinte forma: a
Seção~\ref{sec:fundamentacao} apresenta a fundamentação teórica; a
Seção~\ref{sec:metodo} descreve os materiais e métodos utilizados; a
Seção~\ref{sec:estudo} caracteriza o sistema analisado; a
Seção~\ref{sec:resultados} apresenta os resultados obtidos; a
Seção~\ref{sec:discussao} discute esses resultados; e a
Seção~\ref{sec:conclusao} conclui o trabalho.

\section{Fundamentação Teórica}
\label{sec:fundamentacao}

\subsection{Engenharia reversa de software}

Chikofsky e Cross~\cite{chikofsky1990} definem engenharia reversa como o
processo de análise de um sistema-alvo com o propósito de identificar os
componentes do sistema e seus relacionamentos, criando representações do
sistema em outra forma ou em níveis mais altos de abstração, sem alterar o
próprio sistema. Essa definição distingue engenharia reversa de atividades
correlatas, como reestruturação (\emph{restructuring}) e reengenharia
(\emph{re-engineering}), que envolvem modificação do sistema.

Müller et al.~\cite{muller2000} descrevem a engenharia reversa como parte de
um ciclo mais amplo de compreensão de programas (\emph{program
comprehension}), no qual analistas constroem modelos mentais e artefatos de
documentação a partir de evidências de baixo nível (código-fonte,
configurações, esquemas de banco de dados). Corbi~\cite{corbi1989} já
apontava, ainda na década de 1980, a compreensão de programas como um dos
principais desafios da manutenção de software, dado o custo de reconstruir
conhecimento perdido sobre sistemas legados.

Canfora e Di Penta~\cite{canfora2007} revisam avanços e desafios da área,
destacando a recuperação de arquitetura, a análise de artefatos web e a
geração automática de documentação como fronteiras relevantes de pesquisa,
especialmente à medida que sistemas web e orientados a dados se tornam mais
comuns.

\subsection{Documentação e arquitetura de software}

A arquitetura de software descreve a organização de um sistema em termos de
seus componentes, das relações entre eles e das propriedades observáveis
externamente~\cite{bass2012}. Quando a arquitetura não é documentada
explicitamente, ela ainda existe de forma implícita no código, na
organização de módulos, nas dependências entre eles e nas configurações do
sistema. A \emph{recuperação de arquitetura} (\emph{architecture recovery})
consiste em reconstruir essa estrutura a partir de evidências concretas do
sistema, como estrutura de diretórios, definição de módulos e padrões de
chamadas entre componentes~\cite{koschke2009}.

Sommerville~\cite{sommerville2011} e Pressman e Maxim~\cite{pressman2016}
reforçam que a documentação de requisitos e de arquitetura é um dos ativos
mais valiosos de um projeto de software para fins de manutenção e evolução,
sendo também um dos primeiros artefatos a ficar desatualizado ou ausente em
projetos com restrições de tempo e equipe.

\subsection{UML e modelagem de software}

A \emph{Unified Modeling Language} (UML) é a notação padrão para
representar aspectos estruturais e comportamentais de sistemas orientados a
objetos, incluindo diagramas de casos de uso, classes e
sequência~\cite{omg2017,booch2006}. Fowler~\cite{fowler2005} destaca que
diagramas UML podem ser usados tanto para projetar um sistema quanto para
documentar, de forma simplificada, um sistema já existente, sendo esse
segundo uso o mais próximo do que se aplica neste trabalho: os diagramas
apresentados aqui não representam um projeto planejado previamente, mas uma
reconstrução do comportamento e da estrutura observados no código-fonte.

\subsection{Django e ferramentas de análise}

Django é um framework web em Python que segue o padrão arquitetural
Model-Template-View (variação do Model-View-Controller), oferecendo um
mapeamento objeto-relacional (ORM) integrado, sistema de autenticação, painel
administrativo e um sistema de mensagens \cite{django2026}. O ORM do Django
permite otimizar o número de consultas ao banco de dados por meio dos
métodos \texttt{select\_related} (para relacionamentos de chave estrangeira)
e \texttt{prefetch\_related} (para relacionamentos muitos-para-muitos ou
reversos), mitigando o problema clássico de consultas N+1. O
\emph{Django Debug Toolbar} é uma ferramenta de instrumentação que exibe, em
tempo de desenvolvimento, o número de consultas SQL geradas por requisição, o
tempo de execução e outras métricas de desempenho, sendo amplamente
recomendada como primeiro instrumento de diagnóstico antes de otimizações
mais complexas \cite{django2026}. O banco de dados relacional utilizado no
sistema analisado é o PostgreSQL \cite{postgresql2026}.

\subsection{Trabalhos relacionados}

Os trabalhos de Chikofsky e Cross~\cite{chikofsky1990} e de Müller et
al.~\cite{muller2000} estabelecem os fundamentos conceituais e a taxonomia da
engenharia reversa utilizados neste artigo. Canfora e Di
Penta~\cite{canfora2007} discutem a evolução da área rumo a sistemas mais
complexos e distribuídos, incluindo aplicações web, o que motiva a aplicação
dessas técnicas a um sistema Django. Koschke~\cite{koschke2009} detalha
técnicas de recuperação de arquitetura a partir de artefatos de código,
utilizadas como referência para a reconstrução do diagrama de componentes
apresentado na Seção~\ref{sec:resultados}. Diferentemente desses trabalhos,
que em geral propõem técnicas ou ferramentas de análise automatizada em larga
escala, este artigo tem caráter de estudo de caso, aplicando manualmente um
processo de engenharia reversa assistido por análise estática a um sistema
real de porte pequeno/médio, com foco na síntese de requisitos, regras de
negócio e artefatos UML, e não na criação de uma nova ferramenta de análise.

\section{Materiais e Métodos}
\label{sec:metodo}

O sistema analisado é uma aplicação web desenvolvida em Django 5.2 para
gestão de campeonatos de eSports acadêmicos, contendo sete entidades de
domínio, um módulo de autenticação e um comando de carga de dados de
demonstração. No início deste trabalho, o repositório do sistema não possuía
nenhum documento de requisitos, diagrama ou teste automatizado além do
arquivo padrão gerado pelo Django.

O processo de engenharia reversa foi conduzido em seis etapas, ilustradas na
Figura~\ref{fig:processo}: (1) análise do código-fonte (\emph{models},
\emph{views}, \emph{urls}, configurações e migrações do banco de dados); (2)
extração de evidências (padrões de acesso, filtros de \emph{queryset},
restrições de mixins de autenticação/autorização, campos e relacionamentos de
modelo); (3) inferência de requisitos funcionais, não funcionais e regras de
negócio a partir dessas evidências; (4) modelagem, com a produção de
diagramas de casos de uso, classes, sequência, arquitetura e um modelo
entidade-relacionamento simplificado; (5) análise dos testes automatizados
existentes, incluindo a execução do comando \texttt{manage.py test}; e (6)
análise de desempenho, a partir da inspeção do uso de
\texttt{select\_related}/\texttt{prefetch\_related} nas \emph{views} e da
verificação da configuração do Django Debug Toolbar no projeto.

\begin{figure}[ht]
\centering
\begin{tikzpicture}[
  node distance=0.55cm and 0.35cm,
  every node/.style={font=\small},
  box/.style={draw, rounded corners, align=center, fill=gray!8,
    minimum height=0.9cm, text width=2.05cm, inner sep=2pt},
  >={Stealth[length=2mm]}
]
\node[box] (a) {Análise do\\código};
\node[box, right=of a] (b) {Extração de\\evidências};
\node[box, right=of b] (c) {Inferência de\\requisitos e\\regras};
\node[box, right=of c] (d) {Modelagem\\(UML/DER)};
\node[box, below=1.1cm of c] (e) {Análise de\\testes};
\node[box, left=of e] (f) {Análise de\\desempenho};
\node[box, left=of f] (g) {Consolidação da\\documentação};

\draw[->] (a) -- (b);
\draw[->] (b) -- (c);
\draw[->] (c) -- (d);
\draw[->] (d) -- (e);
\draw[->] (e) -- (f);
\draw[->] (f) -- (g);
\end{tikzpicture}
\caption{Processo de engenharia reversa aplicado ao sistema}
\label{fig:processo}
\end{figure}

Todos os artefatos produzidos (requisitos, regras de negócio, diagramas e
análises) foram fundamentados exclusivamente em evidências localizáveis no
repositório: definições de modelo (\texttt{models.py}), classes de
\emph{view} (\texttt{views.py}), roteamento (\texttt{urls.py}),
configurações (\texttt{settings.py}), migrações de banco de dados e
\emph{templates} HTML. Quando uma informação não pôde ser confirmada por
evidência direta, ela foi explicitamente marcada como não identificada,
seguindo a recomendação metodológica de diferenciar fatos observados de
inferências~\cite{chikofsky1990}. A consolidação final ocorreu em um conjunto
de documentos técnicos estruturados (requisitos, arquitetura/UML, testes,
desempenho e matriz de rastreabilidade), que serviu de base direta para os
resultados apresentados neste artigo.

\section{Estudo de Caso}
\label{sec:estudo}

O sistema analisado -- aqui referido como \emph{Arena eSports} -- tem por
finalidade apoiar a organização de campeonatos de eSports entre múltiplos
campi de uma instituição de ensino, contemplando o cadastro de campi,
modalidades, fases de competição, jogadores, campeonatos, inscrições de
times e jogos (partidas). O sistema também oferece uma tela de listagem
personalizada (``Meus Jogos'') para que um jogador acompanhe as partidas das
quais participa.

O sistema é implementado como uma aplicação Django única (\texttt{website})
sobre um projeto Django (\texttt{pw2026}), utilizando Python 3.12, Django
5.2.13, PostgreSQL como banco de dados relacional (via \texttt{psycopg2}),
\texttt{django-braces} para controle de acesso por grupo,
\texttt{django-crispy-forms} com o tema Bootstrap 5 para formulários, e
\texttt{django-debug-toolbar} para diagnóstico em desenvolvimento. O
front-end utiliza \emph{templates} renderizados no servidor, com Bootstrap 5
para estilização. O deploy é planejado para o Google App Engine, com
Gunicorn como servidor de aplicação.

O modelo de domínio é composto por sete entidades principais: \emph{Campus},
\emph{Modalidade}, \emph{Fase}, \emph{Jogador}, \emph{Campeonato},
\emph{Inscrição} e \emph{Jogo}, além do modelo de usuário nativo do Django.
Essas entidades e seus relacionamentos são detalhados na
Seção~\ref{sec:resultados}.

\section{Resultados}
\label{sec:resultados}

\subsection{Requisitos e regras de negócio}

A Tabela~\ref{tab:rf} apresenta uma seleção dos requisitos funcionais mais
representativos, recuperados a partir das classes de \emph{view} do sistema.
A Tabela~\ref{tab:rn} apresenta as principais regras de negócio
identificadas, em geral implícitas em filtros de \emph{queryset} e em
\emph{mixins} de controle de acesso.

\begin{table}[ht]
\centering
\caption{Requisitos funcionais selecionados}
\label{tab:rf}
\small
\begin{tabular}{@{}p{0.9cm}p{6.6cm}p{4.6cm}@{}}
\toprule
\textbf{ID} & \textbf{Descrição} & \textbf{Evidência} \\
\midrule
RF02 & Cadastro/edição de Modalidades restritos aos grupos \emph{Administrador}/\emph{Organizador}; exclusão restrita a \emph{Administrador} & \texttt{ModalidadeCreate/Update/Delete} \\
RF05 & CRUD de Campeonato, registrando o usuário responsável pelo cadastro & \texttt{CampeonatoCreate.form\_valid} \\
RF08 & Listagem de jogos em que o usuário autenticado participa como jogador de um dos times & \texttt{MeusJogos.get\_queryset} \\
RF10 & Mensagens de sucesso após criação, edição, exclusão e login & \texttt{SuccessMessageMixin}, \texttt{SuccessMessageDeleteMixin} \\
RF11 & Página inicial exibe os 5 campeonatos e as 5 inscrições mais recentes & \texttt{Index.get\_context\_data} \\
\bottomrule
\end{tabular}
\end{table}

\begin{table}[ht]
\centering
\caption{Regras de negócio selecionadas}
\label{tab:rn}
\small
\begin{tabular}{@{}p{0.9cm}p{6.6cm}p{4.6cm}@{}}
\toprule
\textbf{ID} & \textbf{Descrição} & \textbf{Evidência} \\
\midrule
RN01--RN04 & Um usuário só pode editar/excluir os registros de Jogador, Campeonato, Inscrição e Jogo que ele mesmo criou & \texttt{get\_queryset()} sobrescrito em cada \emph{view} de edição/exclusão \\
RN05 & Apenas os grupos \emph{Administrador}/\emph{Organizador} podem gerenciar Modalidades & \texttt{GroupRequiredMixin} \\
RN06 & Exclusões são bloqueadas quando o registro ainda é referenciado por outro & \texttt{on\_delete=PROTECT} na quase totalidade das chaves estrangeiras \\
RN07 & Um jogo relaciona sempre duas inscrições (\texttt{time\_1}, \texttt{time\_2}) e, opcionalmente, um vencedor & Campos de \texttt{Jogo} em \texttt{models.py} \\
\bottomrule
\end{tabular}
\end{table}

Um requisito não funcional recorrente é a exigência de autenticação para a
quase totalidade das operações de escrita (\texttt{LoginRequiredMixin}),
combinada com um único ponto de controle de acesso baseado em grupos
(Modalidade). Também foi identificado que praticamente todas as
\emph{foreign keys} do modelo de dados utilizam \texttt{on\_delete=PROTECT},
uma decisão de projeto que prioriza a preservação do histórico de dados em
detrimento da exclusão em cascata -- um comportamento que só pôde ser
confirmado ao inspecionar as migrações do banco de dados, e não apenas o
código de \emph{views}.

\subsection{Modelagem UML}

A Figura~\ref{fig:casosuso} apresenta o diagrama de casos de uso
simplificado, agrupando os fluxos de CRUD compartilhados por Campus, Fase,
Jogador, Campeonato, Inscrição e Jogo em um único caso de uso, dado que todos
seguem o mesmo padrão de controle de acesso.

\begin{figure}[ht]
\centering
\begin{tikzpicture}[
  every node/.style={font=\footnotesize},
  actor/.style={draw, minimum width=1.6cm, align=center},
  uc/.style={draw, ellipse, align=center, text width=2.5cm, minimum height=0.9cm},
  >={Stealth[length=2mm]}
]
\node[actor] (vis) at (0,0) {Visitante};
\node[actor] (usr) at (0,-2.6) {Usuário\\autenticado};
\node[actor] (org) at (0,-5.2) {Organizador /\\Administrador};

\node[uc] (uc1) at (3.2,0.7) {Consultar página\\inicial};
\node[uc] (uc2) at (3.2,-1.0) {Ver detalhes\\públicos};
\node[uc] (uc3) at (6.4,0.7) {Autenticar-se};
\node[uc] (uc4) at (6.4,-2.0) {Gerenciar Campus,\\Fase, Jogador,\\Campeonato,\\Inscrição, Jogo};
\node[uc] (uc5) at (3.2,-3.6) {Ver Meus\\Jogos};
\node[uc] (uc6) at (6.4,-4.6) {Gerenciar\\Modalidades};

\draw (vis) -- (uc1);
\draw (vis) -- (uc2);
\draw (vis) -- (uc3);
\draw (usr) -- (uc1);
\draw (usr) -- (uc2);
\draw (usr) -- (uc4);
\draw (usr) -- (uc5);
\draw (org) -- (uc4);
\draw (org) -- (uc6);
\end{tikzpicture}
\caption{Diagrama de casos de uso simplificado}
\label{fig:casosuso}
\end{figure}

A Figura~\ref{fig:classes} apresenta o diagrama de classes simplificado,
focado nas sete entidades de domínio e em seus relacionamentos principais,
reconstruído a partir de \texttt{website/models.py} e das migrações do banco
de dados.

\begin{figure}[ht]
\centering
\begin{tikzpicture}[
  every node/.style={font=\tiny},
  cls/.style={draw, align=center, text width=1.7cm, minimum height=0.6cm, fill=gray!8},
  >={Stealth[length=1.6mm]}
]
\node[cls] (user) at (0,3.0) {User\\(auth)};
\node[cls] (campus) at (-3.4,3.0) {Campus};
\node[cls] (jogador) at (-3.4,1.0) {Jogador};
\node[cls] (modalidade) at (3.4,3.0) {Modalidade};
\node[cls] (fase) at (3.4,1.0) {Fase};
\node[cls] (campeonato) at (-1.5,-1.0) {Campeonato};
\node[cls] (inscricao) at (1.5,-1.0) {Inscrição};
\node[cls] (jogo) at (0,-3.2) {Jogo};

\draw[->] (jogador) -- (user);
\draw[->] (jogador) -- (campus);
\draw[->] (campeonato) -- (campus);
\draw[->] (campeonato) -- (user);
\draw[->] (campeonato) -- (modalidade);
\draw[->] (inscricao) -- (campeonato);
\draw[->] (inscricao) -- (modalidade);
\draw[->] (inscricao) -- (user);
\draw[->] (inscricao) -- (jogador);
\draw[->] (jogo) -- (inscricao);
\draw[->] (jogo) -- (fase);
\draw[->] (jogo) -- (modalidade);
\draw[->] (jogo) -- (user);
\end{tikzpicture}
\caption{Diagrama de classes simplificado do domínio}
\label{fig:classes}
\end{figure}

A Figura~\ref{fig:sequencia} apresenta o diagrama de sequência do fluxo de
cadastro de uma Inscrição, um dos casos de uso centrais do sistema,
reconstruído a partir da classe \texttt{InscricaoCreate}.

\begin{figure}[ht]
\centering
\begin{tikzpicture}[
  every node/.style={font=\footnotesize},
  >={Stealth[length=1.8mm]}
]
\foreach \x/\nm in {0/Usuário,2.1/View,4.2/Model,6.0/BD} {
  \node at (\x,3.2) {\nm};
  \draw[dashed] (\x,3.0) -- (\x,0);
}
\draw[->] (0,2.6) -- (2.1,2.6) node[midway,above,font=\scriptsize]{POST inscrição};
\draw[->] (2.1,2.15) -- (4.2,2.15) node[midway,above,font=\scriptsize]{form\_valid()};
\draw[->] (4.2,1.7) -- (6.0,1.7) node[midway,above,font=\scriptsize]{INSERT};
\draw[->] (6.0,1.25) -- (4.2,1.25) node[midway,above,font=\scriptsize]{ok};
\draw[->] (4.2,0.8) -- (2.1,0.8) node[midway,above,font=\scriptsize]{objeto salvo};
\draw[->] (2.1,0.35) -- (0,0.35) node[midway,above,font=\scriptsize]{redirect + toast de sucesso};
\end{tikzpicture}
\caption{Diagrama de sequência: cadastro de Inscrição}
\label{fig:sequencia}
\end{figure}

\subsection{Arquitetura e modelo de dados}

A Figura~\ref{fig:arquitetura} apresenta a arquitetura geral do sistema,
reconstruída a partir da configuração do projeto Django e do arquivo de
implantação (\texttt{app.yaml}).

\begin{figure}[ht]
\centering
\begin{tikzpicture}[
  every node/.style={font=\scriptsize},
  box/.style={draw, rounded corners, align=center, minimum height=0.7cm, text width=2.5cm, fill=gray!8},
  >={Stealth[length=1.8mm]}
]
\node[box] (browser) at (0,3.4) {Navegador\\(Bootstrap 5)};
\node[box] (gunicorn) at (0,2.0) {Gunicorn\\(App Engine)};
\node[box] (urls) at (0,0.7) {URLs / Views\\(Django)};
\node[box] (models) at (-2.8,-0.7) {Models (ORM)};
\node[box] (templates) at (2.8,-0.7) {Templates};
\node[box] (bd) at (-2.8,-2.2) {PostgreSQL};
\node[box] (admin) at (2.8,-2.2) {Admin / Messages\\/ Debug Toolbar};

\draw[->] (browser) -- (gunicorn);
\draw[->] (gunicorn) -- (urls);
\draw[->] (urls) -- (models);
\draw[->] (urls) -- (templates);
\draw[->] (models) -- (bd);
\draw[->] (urls) -- (admin);
\end{tikzpicture}
\caption{Arquitetura geral do sistema}
\label{fig:arquitetura}
\end{figure}

A Figura~\ref{fig:der} apresenta o modelo entidade-relacionamento
simplificado, derivado das migrações do banco de dados.

\begin{figure}[ht]
\centering
\begin{tikzpicture}[
  every node/.style={font=\tiny},
  ent/.style={draw, align=center, text width=1.7cm, minimum height=0.6cm},
  >={Stealth[length=1.6mm]}
]
\node[ent] (user) at (0,3.0) {USER};
\node[ent] (campus) at (-3.4,3.0) {CAMPUS};
\node[ent] (jogador) at (-3.4,1.0) {JOGADOR};
\node[ent] (modalidade) at (3.4,3.0) {MODALIDADE};
\node[ent] (fase) at (3.4,1.0) {FASE};
\node[ent] (campeonato) at (-1.5,-1.0) {CAMPEONATO};
\node[ent] (inscricao) at (1.5,-1.0) {INSCRICAO};
\node[ent] (jogo) at (0,-3.2) {JOGO};

\draw (jogador) -- node[midway,above,font=\tiny]{1:1} (user);
\draw (jogador) -- node[midway,above,font=\tiny]{N:1} (campus);
\draw (campeonato) -- node[midway,left,font=\tiny]{N:1} (campus);
\draw (campeonato) -- node[midway,above,font=\tiny]{N:1} (user);
\draw (campeonato) -- node[midway,right,font=\tiny]{N:N} (modalidade);
\draw (inscricao) -- node[midway,above,font=\tiny]{N:1} (campeonato);
\draw (inscricao) -- node[midway,right,font=\tiny]{N:1} (modalidade);
\draw (inscricao) -- node[midway,above,font=\tiny]{N:1} (user);
\draw (inscricao) -- node[midway,left,font=\tiny]{N:N} (jogador);
\draw (jogo) -- node[midway,left,font=\tiny]{N:1 (x3)} (inscricao);
\draw (jogo) -- node[midway,right,font=\tiny]{N:1} (fase);
\draw (jogo) -- (modalidade);
\draw (jogo) -- (user);
\end{tikzpicture}
\caption{Modelo entidade-relacionamento simplificado}
\label{fig:der}
\end{figure}

\subsection{Testes}

A análise do arquivo \texttt{website/tests.py} revelou que ele contém apenas
o código gerado automaticamente pelo Django (\texttt{from django.test import
TestCase}), sem nenhuma classe ou método de teste implementado. A execução
do comando \texttt{python manage.py test website} confirmou essa ausência,
retornando ``\texttt{Found 0 test(s)}'' e ``\texttt{Ran 0 tests in 0.000s}''.
Não foram encontrados outros diretórios de teste, arquivos
\texttt{test\_*.py}, nem configuração de \emph{frameworks} de teste
adicionais (como \texttt{pytest}) no \texttt{requirements.txt}.

Dessa forma, nenhuma funcionalidade do sistema -- incluindo as regras de
propriedade de registro (RN01--RN04), o controle de acesso por grupo (RN05)
e a integridade referencial garantida por \texttt{on\_delete=PROTECT}
(RN06) -- possui cobertura de teste automatizado. A cobertura de testes
observada é de 0\%, valor obtido diretamente da ausência de testes e da
execução do comando de teste do Django, e não de uma ferramenta de medição
de cobertura (nenhuma foi encontrada configurada no projeto). Essa é
classificada, neste trabalho, como a principal lacuna de qualidade
identificada no sistema.

\subsection{Análise de desempenho}

O projeto declara a dependência \texttt{django-debug-toolbar==6.1.0} no
\texttt{requirements.txt}, com o aplicativo registrado em
\texttt{INSTALLED\_APPS} e \texttt{MIDDLEWARE}, rota habilitada em
\texttt{pw2026/urls.py} e uso restrito a \texttt{INTERNAL\_IPS =
["127.0.0.1"]} em \texttt{settings.py}. Isso indica que a ferramenta está
disponível para diagnóstico em ambiente de desenvolvimento local, mas não
foram encontrados no repositório registros, capturas de tela ou relatórios
com métricas efetivamente coletadas (número de consultas por página, tempo
de resposta). Portanto, a análise de desempenho apresentada nesta seção tem
caráter de \textbf{inspeção estática/diagnóstico preliminar}, e não de
medição experimental.

A Tabela~\ref{tab:desempenho} resume o uso de \texttt{select\_related} e
\texttt{prefetch\_related} identificado nas principais \emph{views} do
sistema, evidenciando uma preocupação deliberada com o problema de consultas
N+1 nas telas de listagem e detalhe mais complexas (Campeonato, Inscrição e
Jogo).

\begin{table}[ht]
\centering
\caption{Otimizações de consulta identificadas por view}
\label{tab:desempenho}
\small
\begin{tabular}{@{}p{2.7cm}p{5.3cm}@{}}
\toprule
\textbf{View} & \textbf{Otimização identificada} \\
\midrule
\texttt{Index} & \texttt{select\_related}+\texttt{prefetch\_related} para campeonatos; \texttt{select\_related} para inscrições recentes \\
\texttt{CampeonatoDetail} & \texttt{select\_related} + \texttt{prefetch\_related} (modalidades e inscrições, via \texttt{Prefetch}) \\
\texttt{InscricaoDetail} & \texttt{select\_related} em cadeia (\texttt{campeonato\_\_campus}) + \texttt{Prefetch} de jogadores \\
\texttt{JogoList}/\texttt{JogoDetail} & \texttt{select\_related} de times, fase, modalidade; \texttt{Prefetch} de jogadores de cada time \\
\texttt{MeusJogos} & \texttt{select\_related} combinado a filtro \texttt{Q()} e \texttt{distinct()} \\
\bottomrule
\end{tabular}
\end{table}

Um ponto de investigação identificado por inspeção de código e de
\emph{templates} é o uso da expressão \texttt{\{\{ object\_list.count \}\}}
nos \emph{templates} de listagem (por exemplo, listagens de Campeonatos e
Inscrições) para exibir o total de registros. Como as \emph{views} envolvidas
já utilizam paginação (\texttt{paginate\_by}), o Django já calcula a
contagem total via o objeto \texttt{paginator}; chamar \texttt{.count()}
novamente sobre o \emph{queryset} no \emph{template} pode gerar uma consulta
\texttt{COUNT} adicional por requisição. Sem uma medição efetiva com o
Django Debug Toolbar, este ponto é reportado apenas como hipótese a validar,
e não como um problema confirmado.

\subsection{Matriz de rastreabilidade}

A Tabela~\ref{tab:rastreabilidade} relaciona alguns dos principais
requisitos aos casos de uso e componentes de implementação correspondentes,
evidenciando a rastreabilidade entre os artefatos recuperados e o código-fonte.

\begin{table}[ht]
\centering
\caption{Matriz de rastreabilidade (seleção)}
\label{tab:rastreabilidade}
\small
\begin{tabular}{@{}p{1.3cm}p{2.6cm}p{3.6cm}@{}}
\toprule
\textbf{Requisito} & \textbf{Caso de uso} & \textbf{Componente/implementação} \\
\midrule
RF05 & Gerenciar Campeonato & \texttt{CampeonatoCreate/Update/Delete} \\
RF08 & Ver Meus Jogos & \texttt{MeusJogos} \\
RN01--RN04 & Gerenciar registros próprios & \texttt{get\_queryset()} sobrescrito \\
RN06 & Excluir registros & \texttt{on\_delete=PROTECT} (models/migrações) \\
\bottomrule
\end{tabular}
\end{table}

\section{Discussão}
\label{sec:discussao}

Os resultados indicam que boa parte da documentação funcional e
arquitetural de um sistema Django de porte pequeno/médio pode ser
reconstruída a partir de evidências estáticas de código, sem acesso aos
desenvolvedores originais. Requisitos funcionais e regras de negócio
ligadas ao controle de acesso (RN01--RN05) foram particularmente fáceis de
recuperar, pois estão concentrados em padrões recorrentes e explícitos no
código das \emph{views} (uso de \texttt{LoginRequiredMixin},
\texttt{GroupRequiredMixin} e sobrescrita de \texttt{get\_queryset()}). Já a
regra de integridade referencial (RN06, uso disseminado de
\texttt{on\_delete=PROTECT}) só pôde ser confirmada de forma completa ao
inspecionar as migrações do banco de dados, e não apenas o código Python de
mais alto nível -- um indício de que a engenharia reversa de sistemas Django
exige a análise conjunta de \emph{models}, \emph{views} e migrações, e não
de apenas uma dessas camadas isoladamente, corroborando a visão de que a
recuperação de arquitetura é um processo multiartefato~\cite{koschke2009}.

Por outro lado, algumas informações se mostraram difíceis ou impossíveis de
recuperar apenas com análise estática. Não foi possível determinar, a partir
do código, como os grupos de autorização \texttt{Administrador} e
\texttt{Organizador} são efetivamente criados em um novo ambiente, nem se
existe uma justificativa de negócio documentada para a ausência de um
mecanismo automático de validação garantindo que o campo \texttt{vencedor}
de um \texttt{Jogo} corresponda a uma das duas inscrições participantes
(\texttt{time\_1} ou \texttt{time\_2}) -- uma regra de negócio que parece
estar implícita no domínio do problema, mas que não é imposta pelo código,
o que ilustra um limite conhecido da engenharia reversa: nem toda intenção
de projeto deixa rastro observável no código-fonte~\cite{chikofsky1990}.

A análise de testes revelou o achado mais crítico deste estudo de caso: a
ausência total de testes automatizados. Esse resultado é relevante para a
discussão teórica porque mostra que a engenharia reversa não se limita a
``redescobrir'' o que o sistema faz, mas também evidencia lacunas de
qualidade -- neste caso, a inexistência de uma rede de segurança para
validar as regras de negócio recuperadas na Seção~\ref{sec:resultados} em
futuras manutenções.

Quanto ao desempenho, a análise estática revelou um uso consistente e
aparentemente deliberado de \texttt{select\_related}/\texttt{prefetch\_related}
nas \emph{views} mais críticas, sugerindo que problemas de consultas N+1 já
foram considerados durante o desenvolvimento. No entanto, a ausência de
métricas reais coletadas com o Django Debug Toolbar limita a conclusão sobre
o desempenho real do sistema em produção: o que se apresenta aqui é uma
análise de risco baseada em código, não uma validação experimental. Essa
distinção -- entre diagnóstico estático e medição experimental -- é
importante para não superestimar a confiabilidade dos achados de desempenho.

Como benefício geral, a engenharia reversa permitiu produzir, em um tempo
relativamente curto, uma base de documentação (requisitos, regras de
negócio, diagramas UML, matriz de rastreabilidade) que não existia antes,
sem necessidade de reescrever ou modificar o sistema. Como limitação, o
processo depende fortemente da qualidade e do estilo do código-fonte
original (nomes de classes e métodos, uso ou não de convenções do
framework); em um sistema com código menos consistente, a taxa de
recuperação de requisitos e regras de negócio tenderia a ser menor do que a
observada neste caso.

\section{Conclusão}
\label{sec:conclusao}

Este trabalho apresentou um estudo de caso da aplicação de engenharia
reversa a um sistema web real desenvolvido em Django que não possuía
documentação formal. A partir da análise estática de \emph{models},
\emph{views}, configurações e migrações, foi possível recuperar requisitos
funcionais e não funcionais, regras de negócio, diagramas UML (casos de uso,
classes e sequência), uma arquitetura de componentes, um modelo
entidade-relacionamento simplificado e uma matriz de rastreabilidade,
respondendo à questão de pesquisa proposta: sistemas Django sem documentação
formal podem ter parte relevante de sua documentação arquitetural e
funcional reconstruída por engenharia reversa baseada em evidências de
código, ainda que com limites claros quando a intenção de projeto não deixa
rastro observável.

A principal contribuição deste trabalho é o relato sistemático desse
processo, incluindo os artefatos recuperados e uma discussão crítica sobre
o que pôde e o que não pôde ser inferido apenas do código-fonte, o que pode
servir de referência metodológica para times que precisem documentar
sistemas legados semelhantes. Como limitações, destacam-se o caráter de
estudo de caso único, a ausência de métricas de desempenho experimentais e a
impossibilidade de validar automaticamente as regras de negócio recuperadas,
dado que o sistema analisado não possui testes automatizados.

Como trabalhos futuros, sugere-se: (i) implementar testes automatizados
priorizando as regras de propriedade de registro e de controle de acesso por
grupo identificadas neste estudo; (ii) validar experimentalmente, com o
Django Debug Toolbar, os pontos de desempenho aqui levantados por inspeção
estática; e (iii) replicar o mesmo processo de engenharia reversa em outros
sistemas Django sem documentação, para avaliar a generalidade dos resultados
obtidos.

\bibliographystyle{plain}
\bibliography{sbc-template}

\end{document}
